\section{\textbf{Introduction}}
\label{sec:introduction}

Intelligence, consciousness, general intelligence, and superintelligence are
frequently placed on a single ascending scale. This treatment obscures the
different questions addressed by these concepts. Intelligence concerns a
system's capacity to use information in selecting among possible actions,
states, or outputs. Consciousness concerns the organization and, in its
phenomenal sense, the subjective character of experience. General intelligence
concerns the breadth, transferability, and adaptive organization of agency.
Superintelligence, we argue, concerns a further capability: moving beyond the
collective human frontier by inventing validated knowledge that humanity did
not previously possess or recognize as an object of inquiry.

Our earlier work defined intelligence at its minimal foundation as
information-dependent selection while distinguishing intelligence from both
functional and phenomenal consciousness \cite{zohourian2026intelligence}. On
this account, intelligence is graded rather than binary. A simple conditional
selection may instantiate an elementary decision structure, whereas advanced
systems can coordinate memory, inference, planning, metacognition, and action.
Increasing sophistication can therefore produce more organized intelligence
without proving the existence or character of subjective experience.

Consciousness should likewise not be restricted in advance to one biological,
humanlike, or computational form. We previously argued that phenomenal,
observable, and functional claims about consciousness must be separated and
that language can serve as a unifying interface among specialized artificial
minds \cite{zohourian2026language}. Language is not merely an input--output
modality under this account. It can provide a shared representational medium
through which perceptual, reasoning, memory, and supervisory systems exchange
information. Such coordination may enable increasingly integrated levels of
functional artificial consciousness without, by itself, establishing
phenomenal experience.

These distinctions become essential when considering artificial general
intelligence (AGI). Our account of AGI as integrated multimodal agency treats
general intelligence and consciousness as related but non-identical dimensions
\cite{zohourian2026agi}. Neither multimodality nor consciousness alone supplies
general intelligence. AGI requires the persistent coordination of perception,
memory, reasoning, objectives, planning, metacognition, and adaptive action
across unfamiliar tasks, domains, and environments. This view is consistent
with accounts that relate intelligence to an agent's ability to achieve goals
across environments rather than to performance on one fixed task
\cite{legg2007universal} and with broader treatments of AGI as adaptable,
domain-general intelligence rather than a collection of isolated task systems
\cite{goertzel2014agi}.

Future work will introduce \emph{AGI-Former}, a neural architecture intended to
operationalize this conception of integrated multimodal agency. The present
paper does not specify that architecture. It instead addresses the conceptual
problem that follows from it: even if a machine realizes AGI, what additional
capability would make it superintelligent?

\subsection{\textbf{From General Intelligence to Superintelligence}}

General intelligence and superintelligence are sometimes treated as
interchangeable descriptions of an increasingly capable system. They instead
identify different functional achievements. AGI denotes integrated,
transferable, and adaptive agency across tasks, domains, modalities, and
environments. Superintelligence denotes capability beyond the collective human
frontier: the capacity of a generally intelligent machine to originate
validated knowledge that humanity did not previously possess or recognize as
an object of inquiry. This retains the established comparative meaning of
superintelligence as capability exceeding human intelligence
\cite{bostrom2014superintelligence}, while specifying the epistemic capability
that distinguishes it from AGI in the present framework.

An AGI could therefore remain bounded by existing human knowledge. It might
reason across disciplines, operate tools, adapt plans, and solve difficult
problems while deriving its objectives and conceptual boundaries from questions
already posed by humans. These abilities would establish general agency, but
they would not alone demonstrate superintelligence. The transition occurs when
the machine ceases to be only a general user of known science and becomes an
originator of new science.

This distinction can be expressed directly:

\begin{equation}
\begin{aligned}
\operatorname{AGI}(x)
&\not\Rightarrow \operatorname{SI}(x),\\
\operatorname{SI}(x)
&\Rightarrow \operatorname{AGI}(x).
\end{aligned}
\label{eq:agi_si_relation}
\end{equation}

The additional condition is not greater speed, scale, memory, or performance
within a predetermined task. It is machine-originated knowledge invention. To
state this condition precisely, the distinction between known unknowns and
unknown unknowns must first be established.

\subsection{\textbf{Known Unknowns and Unknown Unknowns}}

Scientific research commonly begins with a \emph{known unknown}: a question,
target, mechanism, or missing value that researchers have already identified
but not yet resolved. A machine may answer such a question more quickly or
accurately than human researchers and thereby make a major scientific
contribution. The machine is operating beyond prior human performance, but the
epistemic destination was already represented within human inquiry.

An \emph{unknown unknown} is stronger. It is a consequential question,
relation, phenomenon, explanatory possibility, or field of inquiry that
humanity has neither resolved nor recognized as a target. Unknown unknowns are
not simply entries missing from a database. Before they are exposed, humans do
not know that a particular question should be asked, that two bodies of
knowledge should be connected, that an observed anomaly signifies a new
mechanism, or that a new conceptual object is possible.

Let $\mathcal{K}_{H,t}$ denote the collective body of justified human knowledge
at time $t$, and let $\mathcal{Q}_{H,t}$ denote the set of questions, hypotheses,
and research targets explicitly recognized by the human scientific community.
For a candidate epistemic contribution $\kappa$,

\begin{equation}
\begin{aligned}
\kappa \in \mathcal{Q}_{H,t}
\land \kappa \notin \mathcal{K}_{H,t}
&\quad\Longrightarrow\quad
\kappa\ \text{is a known unknown},\\
\kappa \notin \mathcal{Q}_{H,t}
\land \kappa \notin \mathcal{K}_{H,t}
&\quad\Longrightarrow\quad
\kappa\ \text{is an unknown unknown}.
\end{aligned}
\label{eq:known_unknowns}
\end{equation}

The second relation cannot be established solely by observing that a generated
sentence is absent from a training corpus. Semantic novelty, scientific
significance, and epistemic validity are different properties. A system must
identify or construct the previously unrecognized object of inquiry, explain
its relation to existing knowledge, and produce evidence through which the
claim can be independently assessed.

Existing AI-for-science results demonstrate important progress without yet
collapsing this distinction. AI has recovered symbolic relations from data
\cite{udrescu2020aifeynman}, guided mathematicians toward new conjectures and
theorems \cite{davies2021mathematics}, predicted protein structures at high
accuracy \cite{jumper2021alphafold}, and expanded computational catalogues of
candidate stable materials \cite{merchant2023materials}. These achievements
show that machine learning can contribute to discovery at the scientific
frontier. They also remain embedded in research problems, objectives, data,
and validation practices established by human investigators. Under the present
framework, their scientific importance is not in question; what remains to be
established is whether a machine can itself expose an unknown unknown and
originate the knowledge through which it becomes understood.

\subsection{\textbf{Knowledge Invention}}

We use \emph{knowledge invention} for the machine-originated production of a
new epistemic object that lies beyond the collective human frontier and becomes
justified through appropriate validation. The term includes the invention of
new explanatory theories, conceptual structures, scientific methods,
architectures, or research directions. It does not imply that natural laws are
arbitrarily created. A machine may discover a previously unknown phenomenon or
relation, but the representation, explanation, and generalizable body of
knowledge constructed from that discovery constitute an invented epistemic
contribution.

Knowledge fusion and knowledge discovery are principal routes toward knowledge
invention. Knowledge fusion connects bodies of evidence, concepts, models, or
disciplines whose consequential relation was not previously recognized.
Knowledge discovery identifies a previously unknown phenomenon, mechanism,
regularity, entity, or law. Either process may produce knowledge invention when
the machine converts its result into a novel, consequential, and validated
addition to collective knowledge:

\begin{equation}
\begin{aligned}
\operatorname{Fusion}_{x}
&\longrightarrow \kappa_{F},\\
\operatorname{Discovery}_{x}
&\longrightarrow \kappa_{D},\\
\operatorname{Validate}(\kappa_{F}\lor\kappa_{D})=1
&\longrightarrow \operatorname{KI}_{x}(\kappa),
\end{aligned}
\label{eq:routes_to_invention}
\end{equation}

where $\operatorname{KI}_{x}(\kappa)$ denotes knowledge invention attributable
to machine $x$. The arrows represent possible epistemic pathways rather than a
claim that every fusion or discovery yields invention. Most recombinations are
not consequential, many hypotheses are false, and some discoveries answer
questions already posed by humans. Knowledge invention therefore requires a
stronger conjunction:

\begin{equation}
\begin{aligned}
\operatorname{KI}_{x}(\kappa)
\Longleftrightarrow{}&
\kappa \notin \mathcal{K}_{H,t}\\
&\land \kappa \notin \mathcal{Q}_{H,t}\\
&\land \operatorname{Originate}_{x}(\kappa)\\
&\land \operatorname{Consequential}(\kappa)\\
&\land \operatorname{Validate}(\kappa)=1.
\end{aligned}
\label{eq:knowledge_invention}
\end{equation}

The origination condition does not require isolation from humans. Scientific
work is commonly distributed across people, instruments, institutions, and
computational systems. It requires that the machine make the causally essential
move that exposes or constructs the unknown unknown, rather than merely execute
a fully specified human procedure. The consequentiality condition excludes
arbitrary novelty. The validation condition may be satisfied through
controlled experiment, mathematical proof, reproducible computation,
successful engineering realization, or another method appropriate to the
domain. Until validation occurs, the output remains a candidate invention or
hypothesis rather than knowledge.

Recent systems already integrate language models with literature search, code,
planning, laboratory interfaces, and experimental execution
\cite{boiko2023autonomous}, while autonomous laboratories combine computation,
machine learning, robotics, and empirical feedback in closed experimental loops
\cite{szymanski2023autonomous}. Reviews of AI-enabled science similarly show
that learning systems now participate across hypothesis generation, prediction,
experimental design, and analysis \cite{wang2023scientific}. These systems
provide architectural and empirical precursors to machine knowledge invention.
They do not by their existence satisfy Equation~\eqref{eq:knowledge_invention},
because scientific autonomy within a human-specified objective is not identical
to originating an unrecognized object of inquiry.

\subsection{\textbf{The Defining Transition to Superintelligence}}

Under the framework proposed here, AGI is necessary but not sufficient for
superintelligence. An AGI may transfer knowledge, adapt to unfamiliar
environments, and pursue objectives across domains while remaining bounded by
the human knowledge and research targets available to it. Superintelligence
begins when general intelligence can originate and validate knowledge beyond
that collective frontier:

\begin{equation}
\begin{aligned}
\operatorname{SI}(x)
\Longleftrightarrow{}&
\operatorname{AGI}(x)\\
&\land \exists\kappa\;\operatorname{KI}_{x}(\kappa)\\
&\land
\mathcal{K}_{H,t+1}
=\mathcal{K}_{H,t}\cup\{\kappa\},\\
\operatorname{SI}(x)
&\Rightarrow \operatorname{AGI}(x),\\
\operatorname{AGI}(x)
&\not\Rightarrow \operatorname{SI}(x).
\end{aligned}
\label{eq:superintelligence_definition}
\end{equation}

The central thesis of this paper is therefore that \emph{superintelligence is
general intelligence capable of originating validated knowledge beyond the
collective human frontier, including the identification and resolution of
unknown unknowns}. Knowledge fusion and knowledge discovery are not three
parallel requirements alongside knowledge invention. They are principal
generative routes through which knowledge invention may occur.

This definition does not equate superintelligence with omniscience,
infallibility, consciousness, model size, or universal superiority on every
task. It makes a narrower and testable claim: the defining evidence for
superintelligence is a machine-originated epistemic contribution that humanity
did not previously possess or explicitly seek, that materially extends human
understanding, and that survives independent validation. The remainder of the
paper distinguishes this capability from general intelligence,
examines knowledge fusion and discovery as routes to invention, and develops
criteria for attributing knowledge invention to an artificial system without
mistaking recombination, hallucination, or unverified novelty for science.
